\chapter{犯罪客体}

% Source PDF physical page 34; see audit/source_coverage.json.
\begin{EditionOpeningPanel}
重要考点有：保护法益的解释功能，保护法益与保障人权的紧张关系。
\end{EditionOpeningPanel}

[问题]影片《肖申克的救赎》中，被冤枉而坐牢的男主角甲，越狱后如果被抓了回来，要不要判他脱逃罪？（越狱案）

\subsection*{一、犯罪客体}

犯罪客体，是指法律所保护的利益，简称法益。例如，刑法设立故意杀人罪，保护的法益是人的生命。一个罪名的保护法益也被称为保护客体或犯罪客体。犯罪是指侵害法益的行为。（2015年第9题）\footnote{关于犯罪的本质，学术界有“法益侵害说”与“规范违反说”的争论。最早的“规范违反说”认为，犯罪是违反伦理道德规范的行为。但是，这种观点由于混淆了法律与道德，因此在1980年代便被槟弃。对这种“伦理道德规范违反说”，张明楷教授有专门批判，参见张明楷：《行为无价值论与结果无价值论》，北京大学出版社2012年版，第25、26页。}

[注意1]犯罪对象是犯罪客体的载体。犯罪客体是犯罪对象蕴含的利益。例如，财物是盗窃罪的犯罪对象，财物蕴含的财产权是盗窃罪的犯罪客体。

[注意2]法益的解释功能。第一讲讲到目的解释。刑法的保护目的就是法益。目的解释，就是用法益来解释罪名的构成要件。

例如（2022年主观题），乙常年在国外，国内的房子常年是个空房子。甲潜入该房子，住了一个月。甲是否构成非法侵入住宅罪？（1）多数说认为，该罪的保护法益是住宅的居住生活安宁状态。空房子无人居住，不存在居住生活的安宁状态。甲没有侵犯居住生活的安宁状态，因此不构成该罪。（2）少数说认为，该罪的保护法益是居住者的进入许可权。未经乙同意，甲擅闯民宅，侵犯了乙的进入许可权，因此构成本罪。

\subsection*{二、保障人权}

刑法在打击犯罪、保护法益的同时，也应保障人权，严格依据刑法来定罪处罚。刑法不仅是“善良人的大宪章”，也是“犯罪人的大宪章”。过去，我们往往只强调刑法的任务是打击犯罪。实际上，若刑法的任务仅限于此，那么不制定刑法，打击犯罪的效果反倒更好，因为没有约束，想怎么打击就怎么打击。之所以制定刑法，就是为了约束司法机关，在打击犯罪时不能侵犯人权。

% Source PDF physical page 35; see audit/source_coverage.json.
[注意】保护法益与保障人权的冲突。若一味侧重保护法益，则容易认定为有罪。若一味侧重保障人权，则容易认定为无罪。

[结论］两大功能冲突时，优先保障人权。例如，本讲开头“越狱案”中，如果侧重保护法益（监狱秩序），则应判脱逃罪。如果侧重保障人权，则不应判脱逃罪。司法解释的结论是：优先保障人权，不以脱逃罪论处。\footnote{参见1983年8月31日最高人民法院研究室《关于因错判在服刑期“脱逃”后确有犯罪其错判服刑期限可否与后判刑期折抵问题的电话答复》。该司法解释至今有效。男主角甲不构成脱逃罪的理由有两个：一是排除违法事由中的“自救行为”，也即来不及寻求公力救济或公力救济失效时，允许自救行为。二是排除责任事由中的“不具有期待可能性”，也即刑法不能期待被冤枉的人不逃跑。}

